\chapter[动态避障与感知路点方案]{Phase 4 更新：动态避障与感知路点落地方案}
\label{ch:dynamic-perception}

\section{需求来源、编号与实施边界}

本章依据团队于 2026-09-11 转述的上周导师意见新增：机械臂执行任务时将与实验员共处，周围存在化学仪器，需要分别验证进入检测停止和灵活避障。此处 Phase 4 对应项目范围文档中的 Stage 4（端到端受监督演示），拆为 4a 和 4b；不等同于本报告工作包 P4（感知基线）。感知、停止接口和静态碰撞建模必须在固定平台阶段提前开发，不能等到移动集成结束后才开始。

截至本次更新，D-007 已记录 Jazzy 仿真软件基线；D-011 记录的是固定目标、纯机械臂、无工位碰撞物且无附着试管的搬运回归。下面均为\textbf{待实现和待验证方案}，不代表现有系统已具备人员检测、在线避障或视觉定位。相机型号、停止参数及新增 MoveIt 插件尚未选定。

\section{相关研究及其工程用途}

\begin{longtable}{L{3.3cm}L{5.3cm}L{5.9cm}}
\caption{本次新增研究与实现资料；检索日期 2026-09-11}\\
\toprule
来源 & 研究或接口贡献 & ATOM 采用方式与限制 \\
\midrule
\endfirsthead
\toprule
来源 & 研究或接口贡献 & ATOM 采用方式与限制 \\
\midrule
\endhead
Marvel 与 Norcross，2017（在线 2016）\cite{marvel2017ssm} & 分析速度与间距监控中的反应时间、制动距离和测量不确定性 & 4a 的停止距离预算依据；论文不提供本机阈值或合规证明。 \\
Svarny 等，IROS 2019\cite{svarny2019hri} & 比较区域监控与 RGB-D 人体关键点距离监控 & 比较保守占用区与人体细化表示；骨架漏检不能放行未知物体，作者的运行策略不可直接照搬。 \\
Zhu 等，TIE 2024（在线 2023）\cite{zhu2024dynamic} & 级联非线性 MPC 与人工势场组合，预测动态障碍并约束跟踪 & 4b 进阶对照；模型辨识、控制接口和实时计算成本使其不适合直接替代当前基线。 \\
MoveIt Hybrid Planning\cite{moveithybrid2026} & 全局规划器与在线局部规划器由事件逻辑协调 & 作为 4b 候选集成架构；需要自定义重规划、约束及失败逻辑，不是启用即可避人。 \\
MoveIt Servo\cite{moveitservo2026} & 接收关节/末端指令，可做碰撞接近降速与奇异性检查 & 可用于局部执行或视觉伺服；本身不生成绕障路线，也不是安全额定控制器。 \\
AprilTag 3 与 OpenCV PnP\cite{krogius2019apriltag3,apriltagrepo2026,opencvpose2026} & 已知标记尺寸及相机内参时，可从单目图像估计相对位姿 & 工位/试管架定位基线；标记可见不等于试管存在或已经夹稳。 \\
Sajjan 等，ClearGrasp，ICRA 2020\cite{sajjan2020cleargrasp} & RGB-D 透明物体几何重建，处理原始深度失真 & 说明透明试管不能直接依赖深度；深度补全仅在确定性几何方案失败后评估。 \\
\bottomrule
\end{longtable}

前两项支撑检测与停止设计，MPC 研究支撑动态反应候选方案；AprilTag 和 ClearGrasp 分别说明结构化定位的可行性和透明物体深度的局限。已有实验室插入、视觉与接触文献索引见 \texttt{docs/RELATED\_WORK.md}。本次使用论文摘要、作者/机构记录及官方接口文档，不宣称复现了论文性能。MoveIt 在线资料为 Rolling 文档，只用于架构研究；落地时必须检查 D-007 锁定版本并在决策日志记录兼容性试验。

\section{Phase 4a：有人或活动物体进入即停止}

\subsection{最小可实现链路}

推荐以固定在工位上方或斜上方的 RGB-D 相机进行三维占用检测。先验证单相机覆盖；仪器背面、机械臂后侧或其他可进入盲区无法覆盖时，增加不同视角或限制可进入方向，不能以单相机覆盖整个工作区为默认假设。

\begin{enumerate}
  \item 测量并建立桌面、仪器、盖子活动范围、管线和禁入区的静态碰撞模型；透明、反光或深度缺失的仪器仍使用实测保守几何体。
  \item 深度数据按曝光时间变换至机械臂基座坐标，使用当前关节状态做保守机器人自体过滤，再提取未知占用。过滤边界须测试，避免同时删去贴近手臂的手或物体。
  \item 对受监控工作体积及其外侧停止缓冲区检查占用。4a 不要求先识别出“人”：手、推车、其他活动物体均应触发；物体进入后静止也不能被背景更新吸收后解除停止。
  \item 独立于任务规划循环的停止监督模块锁存事件，禁止新动作，向已验证的控制接口请求受控停止，检查停止确认及实测关节速度。取消 action 的返回值不能证明机械臂已停。
  \item 相机断流、数据过期、TF 缺失、关键区域失去可见性或停止确认超时，同样进入禁止运动/故障状态。watchdog 和硬件链路的独立性必须实测。
\end{enumerate}

建议状态为 \texttt{READY -> RUNNING -> STOPPING -> STOPPED -> REVALIDATE}；传感器或执行异常进入 \texttt{FAULT}。进入停止后保持夹持，不能自动松开试管；保持力、失电行为及样品承接方案要由夹爪实测确认。区域清空且持续可观测、定位和夹持状态重新有效、人工确认恢复后，才从当前实测状态重新规划。硬件急停复位不直接恢复旧轨迹。

\subsection{“立即停”的物理含义与距离预算}

不能等人触碰到名义工作范围边缘才忽略制动继续距离。受监控区域应覆盖全部连杆、夹爪、所持试管及任务运动包络，并在外侧增加停止缓冲。借鉴 SSM 的分项方法\cite{marvel2017ssm}，初步工程预算可写为：
\begin{equation}
S_{\mathrm{budget}} = v_H(T_R+T_S) + v_R T_R + B_R + C + Z_H + Z_R .
\end{equation}
其中 $v_H,v_R$ 为保守接近速度上界（m/s），$T_R$ 为感知、处理、通信到制动开始的总反应时间（s），$T_S$ 为后续制动时间（s），$B_R$ 为机器人制动过程的接近距离（m），$C$ 为进入/检测几何余量（m），$Z_H,Z_R$ 为感知和机器人位置误差界（m）。这是\textbf{方案设计预算，不是已验证安全公式或固定阈值}；需考虑各连杆危险点及持物，而不是只用 TCP 速度。

应分姿态、速度、负载和停止接口测量时间及停止扫掠空间；感知端记录曝光时间、处理结束、命令发出、制动开始和静止时刻，并量化时间同步误差。报告最大观测值及统计分布，另加有依据的保守余量；不能仅用平均值或 P95 作为安全上界。普通 RGB-D/ROS 停止原型不能替代经验证的保护装置和硬件急停；人员共处前的硬件保护及风险审查仍遵循项目范围文档。

\section{Phase 4b：受约束的灵活绕行}

\subsection{从停止重规划过渡到在线反应}

4b 保留 4a 的停止监督作为高优先级后备。4a 中有人进入整个受监控工作区即停；4b 只有在新运行模式单独通过验证后，才允许远处障碍进入外层反应区时绕行，而内层停止边界始终触发停止。不能一边沿用“全区进入必停”，一边声称机器人可在同一区域内持续避人。

\begin{enumerate}
  \item \textbf{4b 起步对照：停止、重观测、重规划。}更新 PlanningScene，加入带误差膨胀的动态障碍体，从当前关节状态寻找可行替代路径。执行前再次检查场景版本、时间有效性和停止距离。此项能验证换路，但必须标为间歇式重规划，不能当作连续避障完成。
  \item \textbf{4b 完整目标：运动中更新并反应。}在外层反应区跟踪障碍位置、速度及不确定性，以短时预测包络检查轨迹；评估 Hybrid Planning 的全局/局部组合，必要时用 Servo 作为执行接口。明确规划期限、轨迹替换连续性、控制权仲裁和失效停止；有移动物体的场景快照本身不构成动态安全保证。
  \item \textbf{进阶研究：}基线实时性或平滑性不足时，对照 Zhu 等的 MPC 路线\cite{zhu2024dynamic}。仅在验证位置/速度接口和动力学模型后引入外环预测控制，不假定厂商开放力矩控制。
\end{enumerate}

\subsection{右侧有人时如何决策}

“向左移”是候选路径偏好，不是固定速度命令。规划器检查左侧的仪器、肘部扫掠、自碰撞、关节限位、奇异性及所持试管；完整六自由度 TCP 约束下，六轴机械臂通常没有可任意利用的冗余，不能假定保持末端不动时肘部总能绕开。

若左侧路径满足间距和任务约束，可选择左侧绕行再接回目标；左侧有仪器或通道太窄则停止等待。仅在提管后自由空间搬运阶段允许绕行，同时限制试管倾角、速度和加速度，考虑液体洒出。插管、接触开盖、按键阶段受到器械约束，禁止突然横向躲避：先按已验证停止策略停下，必要时采用单独验证的轴向撤离动作。规划超时、预测失效、跟踪目标丢失、两侧封堵或试管夹持异常均回到停止状态。移动底盘在本阶段保持驻停，底盘和机械臂同时避障须另行验证。

\section{感知选型：RGB 还是 RGB-D}

\begin{table}[H]
\centering
\caption{按任务分配感知，而不是按相机名称判断精度}
\begin{tabularx}{\textwidth}{L{2.3cm}YYY}
\toprule
方案 & 工位/槽位定位 & 人和活动物体 & 推荐角色 \\
\midrule
单目 RGB & 已知尺寸标记 + 内参可恢复公制位姿；未知单点像素不能直接变成三维点 & 仅靠二维检测难以可靠给出任意物体三维间距 & 结构化定位成本基线 \\
RGB-D & RGB 做标记位姿；有效深度可辅助几何核验 & 可构建三维占用，但受遮挡、量程、反光影响 & 同时研究定位和 4a 的首选原型类型 \\
双目 RGB & 通过标定和匹配估计深度 & 可测三维但受纹理与匹配质量影响 & 已有双目资产时对照 \\
腕部相机 & 接近目标时提高图像尺度、支持精对准 & 视角随臂变化，不能独立监控全工作区 & 固定视角误差或遮挡不足时追加 \\
\bottomrule
\end{tabularx}
\end{table}

\decisionbox{建议试验配置，尚非采购决定}{先借用一台 RGB-D，固定在工位斜上方；RGB 通道识别试管架/仪器的 AprilTag 板，深度通道检测三维占用和不透明环境。优先在试管架和仪器刚性基准上设置标记，结合槽位几何定位透明试管。不要把原始深度或学习补全的“玻璃表面”当作精密插入依据。单机是否够用由覆盖试验决定；深度无效区域不能当作空闲。}

选型前用实际安装距离与视场测量标记像素尺寸、检测率、三维有效覆盖、末端定位误差和端到端延迟，再检查近距盲区、快门/同步、照明、线缆和计算负载。只按标称帧率或深度分辨率采购不能证明性能。ClearGrasp 展示透明物体深度补全的研究路线\cite{sajjan2020cleargrasp}，但其训练分布不等于本项目试管、液体、标签和仪器背景。

\section{从相机观测生成试管与仪器 waypoint}

\subsection{路点不是检测框中心}

waypoint 应是带参考坐标系、时间、姿态和任务语义的 TCP 目标。将试管架或仪器定义为 $W$，标记为 $M$，相机光学坐标系为 $C$，机械臂基座为 $B$。约定 ${}^{A}T_D$ 把 $D$ 坐标转换到 $A$。对槽位 $i$、动作阶段 $k$：
\begin{equation}
{}^{B}T_{\mathrm{TCP},i,k}
= {}^{B}T_C\;{}^{C}\hat T_M\;{}^{M}T_W\;{}^{W}T_{\mathrm{TCP},i,k} .
\end{equation}
$ {}^{C}\hat T_M$ 来自 RGB 标记角点与 PnP；$ {}^{M}T_W$ 和工位内目标来自实测治具几何及抓取设计。每个目标分别存储预抓取、夹取、提升、预放置、插入和撤离位姿，以及工具轴、姿态约束和速度限制，不把一个中心点复用于全部动作。盖子用铰链/盖子状态定义动作，按钮用表面法向、行程与接触判据定义，不能照搬试管槽位。

若规划目标链节仍为 D-011 的 \texttt{link\_eef}（记为 $E$），必须使用已标定 TCP 偏移：
\begin{equation}
{}^{B}T_E = {}^{B}T_{\mathrm{TCP}}\left({}^{E}T_{\mathrm{TCP}}\right)^{-1} .
\end{equation}
基座或相机位置改变后须更新标定链；腕部相机使用曝光时的 $ {}^{B}T_E(t){}^{E}T_C$。固定在实验室的相机在底盘移动后不能继续使用旧的常量 $ {}^{B}T_C$；可通过额外可见基座标记或重新注册建立当前相对变换。每条 TF 仅保留一个发布者。

\subsection{实施步骤与拒绝条件}

\begin{enumerate}
  \item 测量试管直径、槽位间距、露出高度、可夹持段、插入口直径/深度和器械活动范围；确认允许安装刚性标记板。若槽中试管可倾斜或高度变化，槽位中心不能代表实际抓取姿态，应增加管口/不透明盖或轮廓观测。
  \item 固定成像分辨率和焦距，标定内参与畸变；标定眼到手或手眼变换，独立测量 TCP、标记到工位及槽位偏置。采用不同位置和朝向，留出未用于拟合的姿态检查最终误差；重投影残差不能替代物理精度。
  \item 检测标记并验证 ID、真实边长、重投影误差和多帧/多标记一致性；平面位姿歧义或遮挡时拒绝输出有效目标，不能沿用无限期旧值。
  \item 用上述变换生成目标；目标消息至少包含时间、坐标系、对象/槽位 ID、位姿协方差、有效状态和标定版本。对接现有目标发布与动作触发分离结构时增加有效性门禁，不把更新 topic 当作自动运动命令。
  \item 执行前核验槽位是否有管、目标工位状态、IK 可达性、姿态限制及完整碰撞模型；夹取后用夹爪/视觉反馈确认，并把试管添加为附着碰撞物。目标已改变则重新观测和规划。
  \item 最终插入使用低速受约束路径；误差预算不足时先增加机械导向或局部视觉，再按测量需要增加力/触觉反馈。无有效位姿、超出标定范围或占用状态不明则禁止开始动作。
\end{enumerate}

\subsection{定位误差与相机比较试验}

平移误差报告 mm，旋转误差报告 deg 或 rad，统一注明单位。误差来源至少包括内参、标记、外参、TCP、槽位几何、管在槽内游隙和抓取偏移；相关误差不能未经分析按独立噪声平方和合并。若径向可用间隙 $c_r=(D_{\mathrm{hole}}-D_{\mathrm{tube}})/2$，可用保守条件 $e_{xy}+L\sin(e_\theta)+m<c_r$ 检查初始可行性，其中 $L$ 为有效插入长度，$m$ 为制造与建模余量，长度须使用同一单位。它不能替代接触与插入验证，也不能把孔径差直接当作径向容差。

使用同一试管架、独立测量基准、相同扰动和留出姿态，比较固定 YAML、RGB 标记、RGB-D 标记加有效深度核验；仅在必要时增加直接试管感知。记录空/有管、偏移/倾斜、标签/液体状态、照明和遮挡。报告位姿偏差、重复性、P95/最大误差、误接受率、拒绝率、定位耗时和完整抓放/插入成功率。

\section{实施顺序与退出证据}

\begingroup
\small
\begin{longtable}{L{2.4cm}L{5.7cm}L{6.4cm}}
\caption{新增工作包与验收门；均为计划}\\
\toprule
顺序 & 交付物 & 退出证据 \\
\midrule
\endfirsthead
\toprule
顺序 & 交付物 & 退出证据 \\
\midrule
\endhead
前置测量 & 传感覆盖图、停止接口表、仪器碰撞模型、TCP/槽位测量表 & 所有进入路径已覆盖或限制；误差和停止参数有测量计划。 \\
P4 感知细化 & 标定配置、工位相对路点生成、有效性门禁与留出数据 & 扰动后定位和操作达到实测容差；错误/过期观测均被拒绝。 \\
Phase 4a & 占用检测、停止监督、锁存恢复、事件日志 & 可控假人/物体进入、静止、遮挡、断流、TF 错误和命令超时均触发规定状态；实测停止扫掠不侵入预设边界。 \\
Phase 4b 起步 & 停止后的替代路径规划与执行前再核验 & 右侧阻挡左侧可行时换路；双侧封堵时等待；记录绕行与等待耗时。 \\
Phase 4b 完整 & 动态跟踪、预测包络、在线路径更新、停止后备 & 外层障碍运动时能更新路径；内层入侵及计算超时仍停止；任务约束、间距和控制连续性通过验证。 \\
\bottomrule
\end{longtable}
\endgroup

建议起始试验规模（不是已有数据或最终可靠性证明）：感知至少 30 个独立测试姿态，每姿态重复 3 次并跨两次重新安装/标定会话；避障采用左侧、右侧、双侧封堵、接触阶段入侵、遮挡/断流五类条件，仿真每类 100 个固定种子，隔离实机用可控替代物每类先做 10 次工程调试。配置速度、负载、入侵轨迹和照明；后续可靠性样本量依据风险和目标失败率另定。

4a 输出端到端停止时间（ms）、停止位移（mm）、最小实测间距（mm）、漏检次数/入侵次数、误停次数/运行小时、失效到停止时间及恢复结果。4b 额外输出更新周期和规划耗时（ms）、截止期违约数、绕行成功/尝试次数、等待时间（s）、试管倾角和任务约束违反次数。使用独立测量或同步外部视频核验停止/间距，报告测量分辨率、同步误差、样本数、极值和置信区间。无人替代物验证先行，再按风险审查结果开展受监督人员试验。
